% ==================================================================
% id:       app-006
% topic:    שימושי הנגזרת
% subtopic: בעיות קיצון
% solution: full
% ==================================================================
עבור הפונקציה $\ds{f(x) =\case{cc}{1 + 4x - x^2, & x \leq 3 \\ \frac{x + 5}{2}, & x > 3}}$ מצאו את הערך הגדול ביותר ואת הערך הקטן ביותר בקטע \([0, 4]\). שרטטו סקיצה של הגרף של \(f(x)\) על מנת לבדוק את תשובתכם.
\hint{נקודת התפר היא מועמדת נוספת לקיצון, לצד קצוות הקטע והנקודות שבהן הנגזרת מתאפסת. טפלו בכל אחד מהקטעים שמשני צדיה בנפרד.}
\begin{solbox}
נשתמש כאן בתוצאה הבאה: אם
$f$
רציפה בקטע
$[a,b]$
וגזירה בקטע
$(a,b)$
אז היא מקבלת את המקסימום ואת המינימום שלה או בקצוות התחום
(כלומר, ב-
$x=a$
או
$x=b$)
או בנקודות שמאפסות את נגזרתה.
אז על מנת להשתמש בתוצאה זו, עלינו לברר האם
$f$
רציפה ב-
$[0,4]$
וגזירה ב-
$(0,4)$.
נסמן
$$.g(x)=1+4x-x^2,\ h(x)=\frac{x+5}{2}$$
נשים לב כי
$g$
ו-
$h$
רציפות בכל הציר הממשי. לכן לרציפותה של
$f$
די לבדוק את נקודת התפר, דהיינו
$x=3$.
מתקיים
$$,\lim_{x\to3^-}f(x)=\lim_{x\to3^-}g(x)=1+4\cdot3-3^2=4=f(3)$$
כמוכן
$$.\lim_{x\to3^+}f(x)=\lim_{x\to3^+}h(x)=\frac{3+5}{2}=4$$
נסיק כי
$f$
רציפה ב-
$[0,4]$.
נשים לב כי
$f$
לא גזירה ב-
$x=3$.
אכן
$$.f'_{-}(3)=-2\neq\frac{1}{2}=f'_{+}(3)$$

לכן, "נפעיל" את התוצאה שהוזכרה בתחילת הפתרון על הפונקציה הנתונה בשני קטעים:
$[0,3]$
ו-
$[3,4]$.
נדון תחילה בקטע
$[0,3]$.
נחשב את ערכי
$f$
בקצוות התחום:
$$.f(0)=1,\ f(3)=4$$

נשאר לבדוק מה הם הערכים של
$f$
בנקודות בהן מתאפסת הנגזרת שלה.
$g'(x)=4-2x=0$
אמ"מ
$x=2$.
מתקיים
$$.f(2)=5$$
עד כה, ערכה המירבי של
$f$
בקטע
$[0,3]$
הוא
$f(2)=5$
וערכה המזערי של
$f$
בקטע
$[0,3]$
הוא
$f(0)=1$.
\end{solbox}

\begin{solbox}
\textbf{המשך פתרון}
נעבור לדיון בקטע
$[3,4]$.
נחשב את הערכים של
$f$
בקצוות התחום:
$$.f(3)=4,\ f(4)=\frac{9}{2}$$
אנו רואים כי
$h'(x)\neq 0$
לכל
$x\in\mathbb{R}$.
לכן ערכה המירבי של
$f$
בקטע
$[3,4]$
הוא
$f(4)=\frac{9}{2}$
וערכה המזערי של
$f$
בקטע
$[3,4]$
הוא
$f(3)=4$.

לאחר השוואת של כל הערכים שהתקבלו עד כה
(בשני הקטעים),
נקבל כי
ערכה המירבי של
$f$
בקטע
$[0,4]$
הוא
$f(2)=5$
וערכה המזערי של
$f$
בקטע
$[0,4]$
הוא
$f(0)=1$.
\end{solbox}
